\section{Signed powers and quotient absorption: a 34-integer sextic}
The incoming quartic41 system uses both a global square guard
$X=2u_0^2$ and a scale coordinate $w_0$ satisfying $X=4qw_0$.
A new bootstrap allows us to replace these by the single cubic equation
\[
 X=4qu^2.
\]
Unlike an ordinary positivity guard, this permits a negative global power
at the start. Its sign is recovered only after its absolute value is
proved to be a power of two. Simultaneously replace the transport
coordinate $z$ by $Z$ and use
\[
 (K_0q+X)(q-F-\alpha)-qF-Z(q-1)=0.
\]
Since $q\mid X$ and $\gcd(q,q-1)=1$, this equation itself forces $q\mid Z$.
Thus $z=Z/q$ is recovered without a separate quotient coordinate.
The intermediate system has 40 integer variables and 25 residuals.

\paragraph{The new negative-power lemma.}
All named coefficient ports are the complete source34 tuple, with
$q=\beta J+1$ and $\beta=2^{1064921}-1$.
The unchanged product guards give $Q=|q|\ge16$,
$Q(Q^2-1)<R<Q^4-2$, and $\eta,\zeta$ the same nonzero sign.
Put $Y=(2v+1)q^3$, $E=XY$, $a=(X+1)Y$, $A=a+2$,
$H=4a+3$, and $k=\eta+\zeta$.
The main norm modulo $H$ excludes $u=0$; hence
$|X|\ge4Q$ and $|E|>3R+1$.

Let $B=|A|$, $P_0=|2XY^2+1|$, $c_0=|c|$, and $k_0=|k|$.
The Pell classifications give
\[
 k_0=2\psi_{P_0}(n),\qquad c_0=\psi_B(p),\qquad p>n\ge1.
\]
The inequality follows from $P_0>B$ and $c_0>k_0$.
Modulo $|E|$, the first Pell recurrence gives
$2n\equiv\epsilon(R+1)$, where
$\epsilon=\operatorname{sgn}(k)$ for $X>0$ and
$\epsilon=\operatorname{sgn}(k)(-1)^{n-1}$ for $X<0$.
Either $n\ge(R+1)/2$, or $n>R$.
These bounds justify the unchanged ordinary strong-rank and signed
auxiliary step-down lemmas, yielding $p=R$.
Consequently the latter branch is impossible and $R=2n-1$.
This argument precedes all assumptions about the sign of $X$.

Suppose $X=-x<0$, and put $Z_0=|Y|$, $r_0=(R-1)/2\ge1$.
Then $x\ge64$, $Z_0\ge4096$, $B\ge(x-1)Z_0-2$, and
$P_0=2xZ_0^2-1$.
Elementary Pell growth gives
\[
 \frac{c_0}{k_0}
 \ge\frac12\left(\frac{(2B-1)^2}{2P_0}\right)^{r_0}
 >\frac{(x-3)^{r_0}}2.
\]
The strict comparison follows from the explicit positive difference
\[
 (2(x-1)Z_0-5)^2-(x-3)(4xZ_0^2-2)
 =4(x+1)Z_0^2-20(x-1)Z_0+25+2(x-3)>0.
\]
Since $c_0/k_0<Z_0+1$, it follows that $Z_0>2^R$ and therefore
$|H|>x2^R+1$ and $|H|>x+2^R$.

For odd $R$ the signed Pell identities modulo $H$ imply, with
$\sigma=\operatorname{sgn}(c)$, one of
\[
 \sigma X\equiv2^R\pmod H,
 \qquad \sigma X2^R\equiv-1\pmod H.
\]
The strict bounds exclude the second and turn the first into equality.
Thus $|X|=2^R$. But $\operatorname{sgn}(q)=\operatorname{sgn}(X)$,
so $q=-2^t$ for some $t\ge0$. This contradicts
$q\equiv1\pmod{2^{1064921}-1}$: the residue of $2^t$ is a positive
power $2^r$, $0\le r<1064921$, never $-1$.
Hence $X,q>0$.
The positive branch now gives $\eta,\zeta>0$ and the inherited
positive-power proof recovers $X=2^R$ and the native compiler endpoint.
A literal quartic41 witness is restored by
$w_0=u^2$, $u_0=2^{(R-1)/2}$, and $z=Z/q$.

Conversely canonical witnesses have $q=2^t$ with $t$ odd,
$R$ odd, and $R>3q$. Thus
$u=2^{(R-t-2)/2}$ is integral and $X=4qu^2$; set $Z=qz$.
This proves the same existential equivalence with finite ZFC refutation.
It uses the inherited compiler and Pell results, not a new kernel
certificate or a proof of consistency.

\paragraph{Low-degree projection.}
Set $\tau_{\rm old}=\tau_{\rm tail}+P-k$ and eliminate
\[
 P=k+E(c-\eta),\quad F_2=f^2,\quad g=y^2-V^2,\quad
 N=q(L+1),\quad C_2=c^2,\quad X=D-ac-\gamma(4a+3).
\]
The first norm becomes
$\tau_{\rm tail}^2+(2\tau_{\rm tail}-k)E(c-\eta)-1$.
All 19 remaining residuals are cubic or lower in 34 unrestricted integer
unknowns. Their sum of squares has exact degree six: its coefficient at
$r^2c^4$ is four, contributed only by $I-(2r+1)c^2$.
The projections have explicit polynomial inverse reconstructions.
This improves the incoming 35-integer sextic by one unknown through a
new signed-power bootstrap, not merely through another monic elimination.
The straightforward octic companion still has 30 unknowns and is not a
new octic variable record.
